\section*{Abstract}
Interpreting and analysing medical images represents a critical yet resource intensive process that often illustrates considerable inter-observer variability among professionals. This project focuses on the development of a system that will use an artificial-intelligence -driven-framework for automated radiology report generation  that combines Retrieval Augmented Generation(RAG) with Large Language Models(LLM).By drawing on contextually relevant validated medical reports and integrating them into the generation process of LLMs ,the system will produce a precise and easy to interpret diagnostic report.This approach reduces the risk of hallucination of LLMs and brings the chances of error to a minimum.It also takes away the risk of fabricated information and strengthens the factual reliability by grounding real clinical data.The implementation of the project covers data preparation ,model training and optimization,system design,architectural integration and extensive evaluation to ensure transparency,dependability ,interpretability, trustworthiness and clinician confidence.Even beyond its technical side,our work will contribute to solving the social problem of  healthcare disparities by providing decision support in underdeveloped areas.It will help to ease the workload of radiologists and extend support to those areas where expert radiologists are hard to find. In conclusion, this RAG-based project  shows the scope of artificial intelligence to augment clinical workflows and can also enhance clinical decision making and promote more equitable access to healthcare services for rural areas. 









\vspace{1cm}
\textbf{Keywords:}Retrieval Augmented Generation(RAG), Large Language Models(LLM),Hallucination Reduction,Clinical Data Grounding,Radiology Report Generation,Rural Healthcare Support,Diagnostic Accuracy,Medical Image Generation,Equitable Access to Healthcare,Healthcare Disparity
\pagebreak
